\documentclass[UTF8,12pt]{ctexart}
\usepackage[a4paper,margin=25mm]{geometry}
% Presentation only. Copy into paper/; contest-specific font requirements prevail.
\RequirePackage{amsmath,amssymb}
% Match the Times-style Latin text used by the CUMCM template.
% XeLaTeX/LuaLaTeX only; other engines retain their existing math fonts.
\RequirePackage{iftex}
\ifPDFTeX\else
  \IfFileExists{unicode-math.sty}{%
    \RequirePackage{unicode-math}%
    \IfFontExistsTF{TeX Gyre Termes Math}{\setmathfont{TeX Gyre Termes Math}}{}%
  }{}%
\fi
\AtBeginDocument{%
  \setlength{\abovedisplayskip}{8pt plus 3pt minus 2pt}%
  \setlength{\belowdisplayskip}{8pt plus 3pt minus 2pt}%
  \setlength{\abovedisplayshortskip}{4pt plus 2pt}%
  \setlength{\belowdisplayshortskip}{6pt plus 2pt minus 2pt}%
}
% One numbered group, aligned relations, optional brace conveys joint constraints.
\newenvironment{modelconstraints}
  {\begin{equation}\left\{\begin{aligned}}
  {\end{aligned}\right.\end{equation}}
\newcommand{\modelunit}[1]{\,\mathrm{#1}}
% Use \numberwithin{equation}{section} only when compatible with the chosen template.

\begin{document}
\section{公式排版示例}
本页仅演示排版。定义、模型和推导分别呈现，符号在首次出现时解释。
令 $x_i$ 表示第 $i$ 类资源的分配量，$c_i$ 表示单位收益，$B$ 表示资源上限。
\begin{equation}
  \max_{\boldsymbol{x}}\quad f(\boldsymbol{x}) = \sum_{i=1}^{n}c_i x_i .
\end{equation}
将必须同时满足的条件组成一个公式组，按关系符对齐：
\begin{modelconstraints}
  \sum_{i=1}^{n}x_i &\le B, \\
  0\le x_i &\le u_i, && i=1,\ldots,n .
\end{modelconstraints}
其中 $u_i$ 为各类资源分配量的上界。分段定义使用专用环境：
\begin{equation}
  g(x)=\begin{cases}
    x^2, & x\ge 0,\\
    0, & x<0.
  \end{cases}
\end{equation}
多步推导按等号对齐，数值代入留到推导之后：
\begin{equation}
\begin{aligned}
  \|\boldsymbol{a}+\boldsymbol{b}\|^2
    &= (\boldsymbol{a}+\boldsymbol{b})^{\mathsf T}(\boldsymbol{a}+\boldsymbol{b})\\
    &= \|\boldsymbol{a}\|^2 + 2\boldsymbol{a}^{\mathsf T}\boldsymbol{b}
       + \|\boldsymbol{b}\|^2 .
\end{aligned}
\end{equation}
公式不进行图片化或整体缩放，编号仅用于需要引用的数学关系。
\end{document}
